\documentclass[11pt]{article}
\usepackage[T1]{fontenc}
\usepackage{lmodern,amsmath,amssymb,amsthm,mathtools,microtype}
\usepackage[margin=1in]{geometry}
\usepackage[colorlinks=true,linkcolor=blue,citecolor=blue,urlcolor=blue]{hyperref}
\newtheorem{theorem}{Theorem}
\newtheorem{lemma}[theorem]{Lemma}
\newtheorem{corollary}[theorem]{Corollary}
\newtheorem{proposition}[theorem]{Proposition}
\theoremstyle{remark}\newtheorem{remark}[theorem]{Remark}
\newcommand{\diag}{\operatorname{diag}}
\newcommand{\Diag}{\operatorname{Diag}}
\newcommand{\one}{\mathbf{1}}
\newcommand{\R}{\mathbb{R}}
\title{Variance reversals and entropy ordering\\in Gaussian variational inference}
\author{Jitendra Prajapati}
\date{25 October 2026}
\begin{document}
\maketitle
\begin{abstract}
Margossian, Pillaud-Vivien and Saul conjectured that increasing the
$\alpha$-divergence parameter above one increases every marginal variance
of the optimal factorized Gaussian approximation to a Gaussian target.
We give a negative answer. For every prescribed $a>1$, an explicit
positive-definite Gaussian target has a unique global optimum whose
marginal variances move in opposite directions as $\alpha$ increases
through $a$. The dimension is at most
$\lceil 8a/(a-1)\rceil+2$; a separate choice has dimension ten at $a=2$.
A rational, dimension-minimal three-dimensional example also has two optimal variances
that decrease by more than five percent between $\alpha=9$ and
$\alpha=10000$. The target is fixed while the divergence parameter
varies. The arguments use exact matrix identities and strict convexity
in precision coordinates, rather than numerical optimization. In contrast to marginal variances,
we also give an elementary proof, using classical matrix tools, that
the optimal Gaussian entropy increases strictly with
$\alpha>1$ for every fixed linear family of precision matrices,
unless the target itself belongs to that family.
\end{abstract}
\noindent\textbf{Scope.} This public research note documents explicit constructions
and proofs; it makes no claim of publication priority. The results concern
an extension conjecture, not the paper's proved ordering theorem on $(0,1)$.

\section{Question and main result}
Fix a non-diagonal positive-definite covariance $\Sigma$ and let
$p=N(0,\Sigma)$. The variational family consists of all Gaussians
$q=N(\nu,D)$ with arbitrary mean $\nu$ and positive diagonal covariance $D$.
We use exactly the convention of \cite{mps}:
\begin{equation}\label{eq:div}
 D_\alpha(p\Vert q)
 =\frac{\int p^\alpha q^{1-\alpha}-1}{\alpha(\alpha-1)},
 \qquad \alpha>1.
\end{equation}
Write $D(\alpha)$ for the covariance of the global optimum.
Definition 2 of \cite{mps} compares divergence objectives by comparing
these diagonal variances componentwise. Their Theorem 3 proves the
$\alpha$-ordering on $(0,1)$. Section 2.1 and the paragraph after
Remark 12 in Section 5.3 conjecture the analogous ordering for
$1<\alpha_1<\alpha_2$.

\begin{theorem}\label{thm:all}
For each prescribed $a>1$, there is an explicit non-diagonal covariance
$\Sigma\succ0$ in dimension
\[
 n=\left\lceil\frac{8a}{a-1}\right\rceil+2
\]
such that the global variational optimum is unique, $D(a)=I$, and
\[
 D'_{00}(a)<0<D'_{11}(a).
\]
Here $\Sigma$ and the dimension are fixed before differentiating the
parameter $\alpha$. There is also such an example in dimension ten
at $a=2$.
\end{theorem}

\begin{corollary}\label{cor:interval}
No nonempty open interval contained in $(1,\infty)$ has universal
componentwise monotonicity of optimal diagonal variances over all
positive-definite Gaussian targets and dimensions.
\end{corollary}
\begin{proof}
Choose $a$ in that interval and the target from Theorem \ref{thm:all}.
For all sufficiently small $h>0$, both $a$ and $a+h$ lie in the interval,
but $D_{00}(a+h)<D_{00}(a)$ and $D_{11}(a+h)>D_{11}(a)$.
\end{proof}

\begin{theorem}\label{thm:three}
For the fixed rational target
\begin{equation}\label{eq:sigma3}
 \Sigma=\frac1{12648}
 \begin{pmatrix}
 5673&4650&4650\\
 4650&9854&-1825\\
 4650&-1825&9854
 \end{pmatrix},
\end{equation}
one has $D(9)=I$ and
\[
 \diag D'(9)=\frac1{14808952}(264375,-3209,-3209)^\top.
\]
Moreover,
\[
 \lim_{\alpha\to\infty}D_{22}(\alpha)
 =\lim_{\alpha\to\infty}D_{33}(\alpha)<\frac{19}{20}.
\]
In fact, $D_{22}(\alpha)=D_{33}(\alpha)<19/20$ for every
$\alpha\ge10000$, so the pair $(9,10000)$ already gives a decrease
exceeding five percent.
\end{theorem}

Theorem \ref{thm:all} makes no fixed-dimension claim as $a\downarrow1$,
and does not assert failure for every prescribed pair of distinct
parameters. The finite comparison in Theorem \ref{thm:three} concerns
the same fixed target at both parameters. Neither theorem
asserts decreasing entropy, poor numerical optimization, or a
calibration failure on real data.

\section{Global optimum and an exact derivative formula}
The following derivation supplies the global-optimality step needed for
a counterexample based on stationary equations.

\begin{lemma}\label{lem:global}
For every $\Sigma\succ0$ and $\alpha>1$, \eqref{eq:div} has a unique
global optimum in the factorized Gaussian family. Its mean is zero,
and its diagonal covariance is a smooth function of $\alpha$ locally
at every point of $(1,\infty)$.
\end{lemma}
\begin{proof}
Put $P=D^{-1}$ and
\[
 \Phi_\alpha(P)=\alpha\Sigma^{-1}-(\alpha-1)P,\qquad
 H_\alpha(D)=\alpha D-(\alpha-1)\Sigma.
\]
The Gaussian integral is finite exactly when $\Phi_\alpha(P)\succ0$.
This is equivalent to $H_\alpha(D)\succ0$ by inversion of
positive-definite matrix inequalities. On that domain its
mean-dependent factor is
\[
 \exp\!\left\{\frac{\alpha(\alpha-1)}2
          \nu^\top H_\alpha(D)^{-1}\nu\right\},
\]
which is uniquely minimized at $\nu=0$. Gaussian integration then
reduces the covariance problem to minimizing
\begin{equation}\label{eq:J}
 J_\alpha(P)=-(\alpha-1)\log\det P-\log\det\Phi_\alpha(P)
\end{equation}
over positive diagonal $P$ with $\Phi_\alpha(P)\succ0$.
This is an open convex domain. For a nonzero diagonal perturbation
$E=\Diag(e_i)$ its Hessian is
\[
 d^2J_\alpha(P)[E,E]
 =(\alpha-1)\sum_i\frac{e_i^2}{P_{ii}^2}
  +(\alpha-1)^2\operatorname{tr}
       (\Phi_\alpha^{-1}E\Phi_\alpha^{-1}E)>0.
\]
The trace term is a squared Frobenius norm, so strict convexity follows.

The domain is nonempty by taking $P$ small, and is bounded
coordinatewise since
$P_{ii}<\alpha(\Sigma^{-1})_{ii}/(\alpha-1)$.
At its boundary a precision entry or $\det\Phi_\alpha$ vanishes.
The corresponding negative logarithm diverges while the other term
is bounded below. A unique interior minimum therefore exists.
The positive-definite Hessian and the implicit-function theorem give
a smooth local branch, which uniqueness identifies with the global
optimum. Infeasible Gaussians have infinite divergence and cannot
improve this minimum.
\end{proof}

In covariance coordinates the equivalent objective is
$\alpha\log\det D-\log\det H_\alpha(D)$, up to constants and a positive
factor. Its stationarity equation, also Proposition 8 of \cite{mps}, is
\begin{equation}\label{eq:station}
 \diag H_\alpha(D)^{-1}=\diag D^{-1}.
\end{equation}

\begin{lemma}\label{lem:K}
Let $a>1$ and let $K$ be a correlation matrix satisfying $K\succ I/a$.
The fixed target
\begin{equation}\label{eq:target}
 \Sigma=\frac{aI-K^{-1}}{a-1}
\end{equation}
is positive definite, and $D(a)=I$. If $v=\diag D'(a)$, then
\begin{equation}\label{eq:derivative}
 [a(K\circ K)-I]v
 =\frac{(K\circ K)\one-\one}{a-1}.
\end{equation}
Here $\circ$ is entrywise multiplication.
\end{lemma}
\begin{proof}
Inverse-order monotonicity gives $K^{-1}\prec aI$, so $\Sigma\succ0$.
At $D=I$, $H_a=K^{-1}\succ0$, and \eqref{eq:station} holds because
$\diag K=\one$. Lemma \ref{lem:global} proves global optimality.
Keeping \eqref{eq:target} fixed while varying $\alpha$ gives
\[
 H'_a=I-\Sigma+a\Diag(v),\qquad
 I-\Sigma=\frac{K^{-1}-I}{a-1}.
\]
Differentiate \eqref{eq:station}, use
$(H^{-1})'=-H^{-1}H'H^{-1}$, and take diagonals to get
\eqref{eq:derivative}. Its coefficient matrix is positive definite:
\[
 a(K\circ K)-I=(aK-I)\circ K\succ0
\]
by the Schur product theorem.
\end{proof}

\section{An explicit reversal at every prescribed parameter}
Fix $a>1$ for the remainder of this section and set
\[
 c=\frac{a-1}{a},\qquad M=\lceil8/c\rceil,\qquad \epsilon=\frac1{10}.
\]
Index the coordinates by a weak coordinate $0$, a hub $1$, and $M$
leaves. Define a real symmetric matrix $R$ with diagonal entries one
and off-diagonal entries
\begin{equation}\label{eq:R}
 R_{01}=\frac1{\sqrt{M+1}},\quad
 R_{0j}=0,\quad R_{1j}=\frac1{\sqrt2},\quad
 R_{ij}=\frac12\quad(i\ne j),
\end{equation}
where $i,j$ in the last three expressions are leaves.

The leaf block is $(I+\one\one^\top)/2$, whose inverse is
$2(I-\one\one^\top/(M+1))$. Eliminating that block leaves the
weak/hub Schur complement
\[
 \begin{pmatrix}
 1&1/\sqrt{M+1}\\
 1/\sqrt{M+1}&1/(M+1)
 \end{pmatrix}\succeq0.
\]
Thus $R\succeq0$. Set
\[
 R_\epsilon=(1-\epsilon)R+\epsilon I,\qquad
 K=I/a+cR_\epsilon.
\]
Then $\diag K=\one$ and $K-I/a=cR_\epsilon\succ0$.
Equation \eqref{eq:target} defines a non-diagonal positive-definite
target, fixed from now on. Lemma \ref{lem:K} gives $D(a)=I$.

To determine the derivative, put
\[
 Z=R\circ R-I,\quad d=c(1-\epsilon)^2,\quad
 B=I+dZ=(1-d)I+d(R\circ R)\succ0.
\]
Here $0<d<1$ and $R\circ R\succeq0$. The off-diagonal entries of $K$
are $c(1-\epsilon)R_{ij}$, so \eqref{eq:derivative} simplifies to
\begin{equation}\label{eq:Binv}
 v=\frac1{a^2c}\bigl(\one-B^{-1}\one\bigr).
\end{equation}
Indeed, the coefficient matrix is $acB$, the right side is
$(d/a)Z\one$, and $dB^{-1}Z=I-B^{-1}$.

By leaf symmetry write
$B^{-1}\one=(h,b,\ell,\ldots,\ell)^\top$ and set
$L=1+d(M-1)/4$. The weak, hub and leaf equations are
\[
 h+\frac d{M+1}b=1,\qquad
 \frac d{M+1}h+b+\frac{dM}{2}\ell=1,\qquad
 \frac d2 b+L\ell=1.
\]
Eliminating $h,\ell$ yields
\begin{equation}\label{eq:schur}
 b=\frac{1-d/(M+1)-dM/(2L)}
          {1-d^2/(M+1)^2-d^2M/(4L)}.
\end{equation}
The denominator is the hub Schur complement of $B\succ0$, hence
strictly positive. On the other hand,
\[
 d(M+1)>\frac{81}{100}c\frac8c=\frac{162}{25}>4.
\]
This is equivalent to $dM/(2L)>1$. The numerator in
\eqref{eq:schur} is negative, so $b<0$ and
$h=1-db/(M+1)>1$. Equation \eqref{eq:Binv} gives
$v_0<0<v_1$, proving the main construction.

For the smaller example at $a=2$, retain \eqref{eq:R} but take
$M=8$ and $\epsilon=1/100$. Then
\[
 c=\frac12,\quad d=\frac{9801}{20000},\quad
 d(M+1)=\frac{88209}{20000}>4.
\]
Every preceding step still applies. Its dimension is ten and its
fixed covariance is $\Sigma=2I-K^{-1}$. This completes the proof of
Theorem \ref{thm:all}. At $\alpha=2$, \eqref{eq:div} is one half
of Pearson's $\chi^2(p\Vert q)$ divergence. The conclusion concerns
changes of the optimum as $\alpha$ varies near two.

\section{A rational three-dimensional example}
Take
\[
 K=\begin{pmatrix}
 1&3/4&3/4\\
 3/4&1&19/50\\
 3/4&19/50&1
 \end{pmatrix},\qquad a=9.
\]
Direct cofactors give
\[
 \det K=\frac{1581}{10000},\qquad
 K^{-1}=\frac1{1581}
 \begin{pmatrix}
 8556&-4650&-4650\\
 -4650&4375&1825\\
 -4650&1825&4375
 \end{pmatrix}.
\]
On the antisymmetric leaf direction, $K-I/9$ has eigenvalue
$8/9-19/50=229/450>0$. On the center/symmetric-leaf subspace its
first diagonal entry is $8/9$ and its determinant is
\[
 \frac89\left(\frac89+\frac{19}{50}\right)
 -2\left(\frac34\right)^2=\frac{47}{16200}>0.
\]
Thus $K\succ I/9$. Formula \eqref{eq:target} gives exactly
\eqref{eq:sigma3}, and Lemma \ref{lem:K} gives $D(9)=I$.

By symmetry set $v=(u,w,w)^\top$. Equation \eqref{eq:derivative}
becomes
\[
 8u+\frac{81}{8}w=\frac9{64},\qquad
 \frac{81}{16}u+\frac{23249}{2500}w=\frac{7069}{80000}.
\]
Its determinant is $1851119/80000>0$. Cramer's rule gives
\[
 (u,w,w)=\frac1{14808952}(264375,-3209,-3209),
\]
proving the local assertion of Theorem \ref{thm:three}.
The strict signs persist under sufficiently small perturbations of
this target, by the same implicit-function argument.

\subsection{Known two-dimensional ordering and minimum dimension}
Li and Turner \cite[Appendix B]{li-turner} already give an exact
characterization and monotonicity of the optimal precisions in two
dimensions. Their parameter $\beta$ corresponds here to
$\alpha=1-\beta$. We include a short alternative proof to make the
minimum-dimension conclusion self-contained.
\begin{proposition}\label{prop:two}
For every non-diagonal positive-definite two-dimensional Gaussian target,
both optimal marginal variances increase strictly on $(1,\infty)$.
Consequently three is the smallest dimension allowing a reversal.
\end{proposition}
\begin{proof}
Fix $a>1$ and let $D_0=D(a)$. Apply the fixed coordinate transformation
$D_0^{-1/2}$ to both target and variational family. The divergence is
invariant under this common invertible change of variables, so the
transformed optimum is
$\widetilde D(\alpha)=D_0^{-1/2}D(\alpha)D_0^{-1/2}$ for the fixed
target $\widetilde\Sigma=D_0^{-1/2}\Sigma D_0^{-1/2}$.
At $a$ it is identity, and
\[
 K=[aI-(a-1)\widetilde\Sigma]^{-1}
   =\begin{pmatrix}1&r\\r&1\end{pmatrix}\succ I/a.
\]
Stationarity gives its unit diagonal; the strict matrix inequality follows
from $\widetilde\Sigma\succ0$. Since the target is non-diagonal, $r\ne0$.
The derivative in Lemma \ref{lem:K} is now
$v_i=D'_{ii}(a)/D_{ii}(a)$, because $D_0$ was frozen. Symmetry of the
two-coordinate system gives
\[
 v_1=v_2=\frac{r^2}{(a-1)((a-1)+ar^2)}>0.
\]
Lemma \ref{lem:global} supplies a smooth optimizer locally everywhere;
uniqueness joins these branches on $(1,\infty)$. Integrating the positive
log-variance derivatives proves the assertion for any two parameters.
For a diagonal target, including dimension one, the variational family
contains the target and its covariance is unchanged. Theorem
\ref{thm:three} therefore attains the minimum possible dimension.
\end{proof}

\section{A decrease exceeding five percent}
We first identify the limit of the actual optimum, without asserting
uniform convergence of the logarithmic barrier at its boundary.

\begin{lemma}\label{lem:limit}
For a fixed positive-definite target in any finite dimension,
$D(\alpha)$ converges as $\alpha\to\infty$ to the unique
minimum-determinant positive diagonal majorant of $\Sigma$.
\end{lemma}
\begin{proof}
Let $P_\alpha=D(\alpha)^{-1}$. It uniquely maximizes
\[
 F_\alpha(P)=\log\det P+
 \frac1{\alpha-1}\log\det[\alpha\Sigma^{-1}-(\alpha-1)P].
\]
Its coordinates are uniformly bounded above for $\alpha\ge2$.
The second term is uniformly bounded above by
\[
 \frac{n\log\alpha+\log\det\Sigma^{-1}}{\alpha-1}\longrightarrow0.
\]
For each fixed positive diagonal $Q\prec\Sigma^{-1}$,
$F_\alpha(Q)\to\log\det Q$. Comparing the optimum with one such $Q$
bounds $\log\det P_\alpha$ below; together with the coordinate upper
bounds this excludes zero coordinates at a subsequential limit.
Any limit $P_*$ satisfies $0\prec P_*\preceq\Sigma^{-1}$.
Comparing with every fixed strictly feasible $Q$ gives
$\log\det P_*\ge\log\det Q$. Scaling makes such $Q$ dense in the
positive closed feasible set. Strict concavity of $\log\det P$ on
diagonal precision gives a unique maximizer, so the full family
converges to it. Inversion gives precisely the unique minimizer of
$\log\det D$ among positive diagonal $D\succeq\Sigma$.
\end{proof}

For \eqref{eq:sigma3}, put
\[
 A=\frac{5673}{12648},\quad B_0=\frac{4650}{12648},\quad
 s=\frac{8029}{12648},\quad t=\frac{11679}{12648}.
\]
Uniqueness and leaf symmetry make the limiting covariance
$\Diag(x,y,y)$. Positive-semidefinite dominance of $\Sigma$ is
equivalent to
\[
 y\ge t>s,\qquad x\ge A+\frac{2B_0^2}{y-s}.
\]
At the optimum the second constraint is tight. The derivative of
$y^2[A+2B_0^2/(y-s)]$ has the sign of
\[
 f(y)=A(y-s)^2+B_0^2y-2B_0^2s.
\]
It is strictly increasing for $y>s$. Exact integer arithmetic gives
\[
 12648^3 f(t)
 =5673\cdot3650^2-4650^2\cdot4379
 =2500(30231417-37873971)<0.
\]
The elementary rational comparisons
$A>56/125$, $s<127/200$, and $B_0^2<17/125$, together with
$19/20-2s<0$, imply
\[
 f(19/20)>
 \frac{56}{125}\left(\frac{63}{200}\right)^2
 -\frac{17}{125}\frac8{25}
 =\frac{583}{625000}>0.
\]
Since $t<19/20$, the unique optimal $y$ lies strictly between
$t$ and $19/20$. Lemma \ref{lem:limit} proves the final assertion
of Theorem \ref{thm:three}, including a strict five-percent drop at
some finite $\alpha>9$. The explicit finite comparison is proved next.


\subsection{An explicit finite parameter pair}
For the same rational target, retain $A,B_0,s,t$ above. Fix $a>1$,
put $r=a-1$ and $k=a/r$, and for $y>t/k$ define
\[
 q=ky-s,\qquad d=ky-t,\qquad x=A+\frac{2B_0^2}{q}.
\]
In the center/symmetric-leaf/antisymmetric-leaf basis,
$[a\operatorname{diag}(x,y,y)-r\Sigma]/r$ has a block
\[
 \begin{pmatrix}kx-A&-\sqrt2B_0\\-\sqrt2B_0&q\end{pmatrix}
 \quad\hbox{and the entry }d.
\]
Its block determinant is $(k-1)xq=xq/r>0$, so the covariance is
feasible. The first-coordinate stationarity equation holds identically;
the two leaf equations reduce to
\begin{equation}\label{eq:finite-E}
 E_a(y):=\frac{y}{2r}\left(\frac1q+\frac1d\right)
 +\frac{yB_0^2}{q(Aq+2B_0^2)}-1=0.
\end{equation}
Every root is therefore a feasible stationary point and is the unique
global optimum by Lemma~\ref{lem:global}. The function is continuous on
$(t/k,\infty)$ and tends to positive infinity at its left endpoint.

Set $y_0=19/20$, $q_0=y_0-s$, and $d_0=y_0-t$.
Exact rational comparison gives
\[
 \frac3{10}<q_0<\frac13,\quad d_0>\frac1{40},\quad
 0<A<\frac12,\quad 0<B_0^2<\frac17,
\]
and the estimates already used above give
\[
 f(y_0)>\frac{583}{625000}>\frac9{10000},\qquad
 q_0(Aq_0+2B_0^2)<\frac{19}{126}<\frac16 .
\]
At $y_0$, the actual $q$ and $d$ exceed $q_0$ and $d_0$.
The last rational term in \eqref{eq:finite-E} decreases as $q$ increases.
Using $y_0<1$ therefore yields
\[
 E_a(y_0)
 <\frac{65}{3r}
   -\frac{f(y_0)}{q_0(Aq_0+2B_0^2)}
 <\frac{65}{3r}-\frac{27}{5000}<0
 \quad(a\ge10000).
\]
For the final inequality, $65/29997<1/400<27/5000$.
The intermediate value theorem supplies a root with
$t/k<y<19/20$. Strict convexity identifies its covariance with
$D(a)$, completing the finite assertion of Theorem~\ref{thm:three}.
The accompanying exact-rational verifier checks the numerical matrix
identities and the scalar sign at $a=10000$; the existence and
global-optimality argument is the analytic proof above.

\section{Entropy ordering for linear Gaussian precision families}
Individual variances can reverse, but a universal aggregate ordering
survives. The following theorem also covers linear precision structures
beyond the factorized family. This is an application of classical
log-determinant and projection inequalities, not a claim of a new
fundamental barrier principle; see the discussion below.

\begin{theorem}\label{thm:entropy}
Let $\mathcal L$ be a real linear subspace of the symmetric $n\times n$
matrices containing a positive-definite matrix. Fix $p=N(\mu,\Sigma)$
with $\Sigma\succ0$. Minimize \eqref{eq:div} over
$q=N(\nu,P^{-1})$, with arbitrary $\nu$ and
$P\in\mathcal L$, $P\succ0$.
For every $\alpha>1$ the global optimum is unique and depends smoothly
on $\alpha$. If $\Sigma^{-1}\notin\mathcal L$, then
\[
 \frac{d}{d\alpha}\log\det P(\alpha)^{-1}>0
 \quad\text{for every }\alpha>1.
\]
If $\Sigma^{-1}\in\mathcal L$, the optimum equals $p$ for every $\alpha$
and its entropy is constant.
\end{theorem}

\begin{proof}
Translate the target mean to zero. Gaussian integration as in
Lemma \ref{lem:global} makes the optimal variational mean zero.
Write $\Lambda=\Sigma^{-1}$. The precision objective is again
\[
 J_\alpha(P)=-(\alpha-1)\log\det P-\log\det\Phi_\alpha(P),
 \qquad \Phi_\alpha(P)=\alpha\Lambda-(\alpha-1)P.
\]
Its domain in $\mathcal L$ is $P\succ0$, $\Phi_\alpha(P)\succ0$.
It is open, convex and nonempty: a sufficiently small positive multiple
of any positive-definite member of $\mathcal L$ is feasible.
It is bounded because $0\prec P\prec\alpha\Lambda/(\alpha-1)$.
The objective diverges at every boundary, with either determinant
vanishing and the other negative logarithm bounded below.
For a nonzero $E\in\mathcal L$, its Hessian is
\begin{align*}
 d^2J_\alpha(P)[E,E]
 &=(\alpha-1)\operatorname{tr}(P^{-1}EP^{-1}E)\\
 &\quad+(\alpha-1)^2
   \operatorname{tr}(\Phi_\alpha^{-1}E\Phi_\alpha^{-1}E)>0.
\end{align*}
Thus a unique interior global optimum exists. The restricted Hessian
is invertible, giving smooth local branches by the implicit-function
theorem; uniqueness joins them globally.

Fix $a>1$ and put $P_0=P(a)$. Apply the single fixed coordinate
transformation $x\mapsto P_0^{1/2}x$. The transformed quantities are
\[
 \widetilde P(\alpha)=P_0^{-1/2}P(\alpha)P_0^{-1/2},\quad
 \widetilde\Lambda=P_0^{-1/2}\Lambda P_0^{-1/2},\quad
 \widetilde{\mathcal L}=P_0^{-1/2}\mathcal L P_0^{-1/2}.
\]
The transformed space contains $I$, the optimizer at $a$ is $I$, and
the derivative of the log determinant of covariance is unchanged.
Drop tildes. Let $\Pi$ be Frobenius orthogonal projection onto this
transformed $\mathcal L$, and set
\[
 H=[a\Lambda-(a-1)I]^{-1}\succ0.
\]
The first-order condition is
$\Pi(P^{-1}-\Phi_\alpha^{-1})=0$. At $a$ it gives
\begin{equation}\label{eq:projectstation}
 \Pi H=I,\qquad \operatorname{tr}(HF)=\operatorname{tr}F
 \quad(F\in\mathcal L).
\end{equation}
The projection here uses the transformed subspace; no invariance of
the original Frobenius projection under congruence is assumed.

Let $E=P'(a)\in\mathcal L$. Differentiation, with target and subspace
fixed, gives
\[
 \mathcal A E=\frac1a\Pi(H-H^2),\qquad
 \mathcal A F=\Pi\{F+(a-1)HFH\}.
\]
Indeed $\Phi'_a=\Lambda-I-(a-1)E$ and
$\Lambda-I=(H^{-1}-I)/a$. The operator $\mathcal A$ is self-adjoint
and positive definite on $\mathcal L$, since
\[
 \langle F,\mathcal A F\rangle
 =\|F\|_F^2+(a-1)\|H^{1/2}FH^{1/2}\|_F^2>0
 \quad(F\ne0).
\]
Using \eqref{eq:projectstation} and
$\mathcal A I=I+(a-1)\Pi H^2$, we obtain
\[
 E=\frac{a\mathcal A^{-1}I-I}{a(a-1)}.
\]
At normalized $P=I$, the log-determinant derivative of the covariance
is therefore
\begin{equation}\label{eq:entropyderiv}
 -\operatorname{tr}E
 =\frac{n-a\langle I,\mathcal A^{-1}I\rangle}{a(a-1)}.
\end{equation}

Consider the direct sum of two symmetric-matrix spaces with Frobenius
inner products, and define
\[
 \mathcal T(F)=(F,\sqrt{a-1}\,H^{1/2}FH^{1/2}),\qquad
 z=(I,\sqrt{a-1}\,I).
\]
Then $\mathcal T^*\mathcal T=\mathcal A$. By
\eqref{eq:projectstation},
$\langle z,\mathcal T(F)\rangle=a\operatorname{tr}F$ for all
$F\in\mathcal L$, so $\mathcal T^*z=aI$ and $\|z\|^2=an$.
The projection of $z$ onto $\operatorname{im}\mathcal T$ is
$\mathcal T(a\mathcal A^{-1}I)$ and has squared norm
$a^2\langle I,\mathcal A^{-1}I\rangle$. Consequently
\begin{equation}\label{eq:entropyresidual}
 \frac{d}{d\alpha}\log\det P(\alpha)^{-1}\Big|_{\alpha=a}
 =\frac{\|z-\mathcal T(a\mathcal A^{-1}I)\|^2}{a^2(a-1)}
 \ge0.
\end{equation}
Equality forces $z=\mathcal T(F)$ for some $F\in\mathcal L$.
Its first component gives $F=I$ and its second gives $H=I$,
hence $\Lambda=I$ in normalized coordinates. Undoing the transformation
gives the original $\Sigma^{-1}=P_0\in\mathcal L$.
Conversely, if the target precision belongs to $\mathcal L$, exact
fitting is feasible and is the unique optimum for every $\alpha$.

For a target outside the family, \eqref{eq:entropyresidual} is strictly
positive at every $a>1$. Integration yields strict ordering between any
two parameters. Gaussian differential entropy is
$n\log(2\pi e)/2+\log\det P(\alpha)^{-1}/2$, so the same strictness
holds for entropy.
\end{proof}

Theorem \ref{thm:entropy} includes diagonal and block-diagonal precision,
fixed graphical zero patterns, and symmetric Toeplitz precision
subspaces. It does not assert componentwise or Loewner covariance
ordering, or extend to arbitrary affine, nonlinear or inequality
constraints. The target and the linear subspace must remain fixed.
It also makes no claim for non-Gaussian targets.

\paragraph{Relation to classical barrier inequalities.}
The change of variables $\mu=\alpha-1$ and
$Q=\mu P/(1+\mu)$ makes the feasible set independent of $\mu$.
Up to constants, the precision objective is
$\mu f(Q)+g(Q)$ with
$f(Q)=-\log\det Q$ and $g(Q)=-\log\det(\Sigma^{-1}-Q)$.
At a differentiability point write $d=Q'_\mu$ and
$h=-Df(Q)[d]$. Differentiated optimality gives
\[
 Dg(Q)[d]=\mu h,\qquad
 \mu D^2f(Q)[d,d]+D^2g(Q)[d,d]=h.
\]
The standard barrier gradient inequality
$|DF(Q)[d]|^2\le nD^2F(Q)[d,d]$ for each log-determinant
barrier (see \cite[Definition 2.1, Eq.~(2.4)]{nemirovski-todd}) gives
\[
 \frac{\mu(1+\mu)}{n}h^2\le h,\qquad
 h\le\frac{n}{\mu(1+\mu)}.
\]
Consequently the optimized covariance log determinant
$-\log\det Q_\mu+n\log(\mu/(1+\mu))$ is nondecreasing.
Equality in this entropy derivative bound requires
$h=n/[\mu(1+\mu)]>0$. Equality in both trace Cauchy--Schwarz
inequalities then forces
$Q=\mu\Sigma^{-1}/(1+\mu)$, hence exact fitting.
This is a short alternative classical derivation of
Theorem~\ref{thm:entropy}. The precise prior occurrence of the
Gaussian consequence has not been established here.

\section{Scope and prior-work boundary}
These constructions settle the explicit above-one extension
conjecture in \cite{mps} in the negative. The source's ordering theorem on $(0,1)$,
its forward-KL comparison for $\alpha>1$, its entropy-matching
result, and its impossibility theorem remain unaffected.

The obstruction occurs at unique global optima of smooth
positive-definite Gaussian targets. It cannot be attributed to a
Monte Carlo gradient, local minimum, optimizer tie, or singular target.
It does show that increasing $\alpha>1$ is not a universally
componentwise monotone way to enlarge variational uncertainty.
It does not establish a real-data consequence, an entropy reversal,
or failure of any inference algorithm's convergence guarantee.

The conjecture appears in the official JMLR paper and arXiv v5.
The two-dimensional characterization is prior work
\cite[Appendix B]{li-turner}. The entropy result is an elementary
consequence of classical log-determinant barrier inequalities
\cite{nemirovski-todd}; no foundational novelty is claimed for it.
This release makes no claim of being the first construction or proof
of the higher-dimensional reversals. A literature search and a
timestamp do not establish priority.

Monotonicity of a variational bound at fixed distributions
\cite{li-turner}, decrease along optimization iterations \cite{daudel},
and ordering of unnormalized mass \cite{minka} are different statements
from ordering the entropy of the normalized optimizer as $\alpha$ varies.

\paragraph{Methods and verification.}
AI tools were used extensively in the literature search, derivations,
writing, and checking of this note. The accompanying Python verifier
uses exact rational arithmetic for the explicit finite matrix example
and finite-parameter comparison. It does not formally verify the
analytic theorems. Separate computational checks and AI proof reviews
do not constitute independent human peer review. No numerical
variational optimizer, trained-model experiment, or practical
inference improvement is reported.

\begin{thebibliography}{9}
\bibitem{nemirovski-todd}
A.~S. Nemirovski and M.~J. Todd.
\newblock Interior-point methods for optimization.
\newblock \emph{Acta Numerica}, 17:191--234, 2008.
\newblock \url{https://www2.isye.gatech.edu/~nemirovs/Published.pdf}.

\bibitem{mps}
C.~C. Margossian, L. Pillaud-Vivien, and L.~K. Saul.
\newblock Variational Inference for Uncertainty Quantification:
an Analysis of Trade-offs.
\newblock \emph{Journal of Machine Learning Research},
26(202):1--41, 2025.
\newblock \url{https://jmlr.org/papers/v26/24-0878.html}.
\newblock Retrieved version also available as
\href{https://arxiv.org/abs/2403.13748v5}{arXiv:2403.13748v5}.
\bibitem{li-turner}
Y. Li and R.~E. Turner.
\newblock R\'enyi Divergence Variational Inference.
\newblock \emph{Advances in Neural Information Processing Systems} 29,
2016. Appendix B.
\newblock \url{https://arxiv.org/abs/1602.02311v3}.

\bibitem{daudel}
K. Daudel, R. Douc, and F. Roueff.
\newblock Monotonic Alpha-divergence Minimisation for Variational Inference.
\newblock \emph{Journal of Machine Learning Research}, 24, 2023.
\newblock \url{https://www.jmlr.org/papers/volume24/21-0249/21-0249.pdf}.

\bibitem{minka}
T. Minka.
\newblock Divergence measures and message passing.
\newblock Microsoft Research technical report MSR-TR-2005-173, 2005.
\newblock \url{https://tminka.github.io/papers/message-passing/minka-divergence.pdf}.
\end{thebibliography}
\end{document}
